\section*{ACCIDENTES RADIOLÓGICOS}

Un accidente radiológico puede ser provocado por:

\begin{enumerate}
    \item [a)] Error humano.
    \item [b)] Falla del equipo.
    \item [c)] Falla y/o carencia de equipo de protección radiológica.
    \item [d)] Imprevistos.
\end{enumerate}

La prevención de los accidentes se hace en función de estos cuatro factores, tomando como base:

\begin{enumerate}
    \item [a)] La preparación del personal ocupacionalmente expuesto:

        Los técnicos radiólogos deberán estar conscientes de los riesgos que implica una exposición a la radiación superior al límite anual establecido (sobreexposición), dependiendo de la magnitud de la misma y de si ésta se recibió sólo en ciertas partes del organismo o en todo el cuerpo.

        Deberán conocer el funcionamiento de los equipos de gammagrafía, de los accesorios de los mismos y de los equipos usados en protección radiológica.

        Deberán saber seleccionar, de acuerdo al tipo de trabajo que tengan que desarrollar, los arreglos y/o dispositivos de los gammágrafos, a fin de que reciban la menor exposición posible.

        Deberán no solo conocer, sino realizar, casi en forma automática los procedimientos de revisión, señalización, ensamble y desmontaje del equipo, etc. De ahí la importancia de que el entrenamiento sea teórico y práctico, y que esta última fase esté siempre supervisada por una persona de mayor experiencia. En relación a las acciones de emergencia, debe recordarse que cuando el accidente se presenta, el individuo se encuentra bajo una gran tensión y lo más probable sea que olvide lo que debe hacer, si no se le ha dicho en repetidas ocasiones qué acciones debe seguir. Habrá ocasiones, cuando se sabe en qué lugar se encuentra la fuente y que ninguna persona corre el riesgo de una sobreexposición, en que sea recomendable que el técnico se dé tiempo para calmarse y pensar con detenimiento cómo debe actuar, en lugar de actuar precipitadamente y cometer errores que pueden resultar muy caros.
\end{enumerate}
